\section{Identity: canonical forms and content addressing}\label{sec:identity}

A \VG{} node is named by hashes. The \emph{body hash} covers the op graph, the \emph{contract
hash} the canonical contract, and the \emph{node hash} the format version, the effect set,
the signature and the other two. All are BLAKE3 hashes in derive-key mode, one context string
per kind of hash, so hashes of different kinds are computed by different keyed functions. This
section characterizes exactly when two bodies share a hash, shows that a shared node hash
implies a shared meaning, and explains why semantic equivalence needs a different mechanism.

\subsection{Canonicalization}

The specification canonicalizes a node in seven steps: validate the structure; resolve labels
to hashes and drop identity ops and debug information; hash bottom-up; sort the operands of
commutative ops by hash; merge ops with equal hashes and drop ops on which no output, effect
chain or live invariant depends; number the remaining ops by a post-order depth-first walk
from the outputs, in port order; serialize the sections deterministically. The hash of an op
covers its opcode, mode, attribute word, result types and the hashes of its operands in port
order, after the commutative sort. An input hashes as its position and type, a constant as its
type and exact bytes, a callee as its node hash, and an op result as its op's hash and its
result index.

\subsection{Terms and unraveling}

\begin{definition}[Terms]\label{def:terms}
\emph{Operand terms} $t$ and \emph{op terms} $u$ are given by
\[
t ::= \mathsf{in}(j,\tau) \mid \mathsf{const}(\tau,b) \mid \mathsf{callee}(h) \mid
      \mathsf{inv}(h) \mid \langle u\rangle_s
\qquad\qquad
u ::= f(t_1,\dots,t_n),
\]
where $b$ is a byte string, $h$ a hash, $s$ a result index and $f$ an op symbol of arity $n$.
The \emph{unraveling} $\unr$ of an operand of a well-formed node is defined by recursion on the
op index: $\unr(\mathsf{in}(j))=\mathsf{in}(j,\tau_j)$,
$\unr(\mathsf{const}(c))=\mathsf{const}(\tau_c,\enc(c))$,
$\unr(\mathsf{callee}(h))=\mathsf{callee}(h)$, $\unr(\mathsf{inv}(h))=\mathsf{inv}(h)$, and
$\unr(\mathsf{op}(i,s))=\langle f_i(\unr(a_{i,1}),\dots,\unr(a_{i,n_i}))\rangle_s$.
\end{definition}

Unraveling forgets sharing: a body is a term graph~\cite{ariola1996equational}, and its
unraveling is the tuple of trees it denotes.

\begin{definition}[Equality up to commutativity]\label{def:eqc}
$\eqC$ is the least congruence on terms that contains $f(t_1,t_2)\eqC f(t_2,t_1)$ for every
$f\in\mathit{Comm}$ and $\op{FMA}(t_1,t_2,t_3)\eqC\op{FMA}(t_2,t_1,t_3)$.
\end{definition}

\subsection{Structural encodings and the capacity condition}

Every hash is computed from a byte string, its \emph{preimage}, built from fixed-width
integers, 32-byte hashes and type encodings. A type is encoded as a tag byte, an argument-count
byte, and its arguments, each of which is either a nested type encoding or a 32-bit immediate
(a bit width, an array length, a tensor dimension or a handle kind); the tag determines which.
The preimages hashed under the op context are
\begin{align*}
\pre(\mathsf{in}(j,\tau)) &= \mathtt{00}\cat u_{32}(j)\cat\enc_T(\tau) &
\pre(\mathsf{callee}(h)) &= \mathtt{03}\cat h \\
\pre(\langle u\rangle_s) &= \mathtt{01}\cat\Hop(\pre(u))\cat u_8(s) &
\pre(\mathsf{inv}(h)) &= \mathtt{04}\cat h \\
\pre(\mathsf{const}(\tau,b)) &= \mathtt{02}\cat\enc_T(\tau)\cat u_{32}(|b|)\cat b
\end{align*}
\begin{align*}
\pre(f(t_1,\dots,t_n)) = {}& \mathtt{05}\cat u_{16}(\mathit{opcode})\cat u_8(\mathit{mode})\cat
  u_{32}(\mathit{attr})\cat u_8(r)\cat\enc_T(\rho_1)\cat\dots\cat\enc_T(\rho_r)\\
  & {}\cat u_{32}(n)\cat h_{\pi(1)}\cat\dots\cat h_{\pi(n)},
\end{align*}
where $h_k=\Hop(\pre(t_k))$ and the permutation $\pi$ sorts the commutative positions
($\{1,2\}$ for $f\in\mathit{Comm}$ or $f=\op{FMA}$) by increasing byte string and fixes the
others. We write $\hash(t)=\Hop(\pre(t))$. Ties in the sort are harmless: two operands with the
same hash contribute the same bytes in either order. The hash of an invariant covers the ports
it names through their operand hashes; since an invariant never names a port that depends on the
fold that uses it (\ref{wf:order}), these definitions are well founded. The prototype rejects a
violation with the message ``an INVARIANT names a port that depends on the fold that uses it''.

\begin{assumption}[Field capacity]\label{cond:capacity}
Every integer written into a $w$-bit field is below $2^w$. For the fields above this means: at
most 255 tuple members, at most 254 tensor dimensions (the argument count of a tensor is one
more than its number of dimensions), at most 255 results per op, and result indices below 256.
\end{assumption}

\begin{lemma}[Prefix-free encodings]\label{lem:prefixfree}
Under \cref{cond:capacity}, $\enc_T$ is prefix-free: no type encoding is a proper prefix of
another, and a concatenation of type encodings determines the sequence of types. Preimages of
different kinds differ in their first byte, and equal preimages have equal fields: the same
leaf; or the same op hash and result index; or the same op symbol, arity and sequence of operand
hashes.
\end{lemma}
\begin{proof}
For types, by induction on nesting depth. A decoder reads the tag, which fixes the kind of each
argument (for example, an array has a type and then an immediate, a tuple has only types); it
reads the count, which by \cref{cond:capacity} is the true number of arguments; and it reads
each argument, an immediate of four bytes or, by induction, a uniquely delimited type
encoding. So each encoding is parsed in exactly one way and ends at a determined position,
which excludes proper prefixes. For preimages, the first byte fixes the kind; the remaining
fields have fixed widths, are 32-byte hashes, or are type encodings, which are uniquely
delimited, and counts ($r$, $n$, $|b|$) are exact by \cref{cond:capacity}. So equal preimages
have equal fields.
\end{proof}

\begin{remark}[Why the condition is not a formality]\label{rem:capacity}
Without \cref{cond:capacity} the encoding is ambiguous. If the member count of a tuple is
written as one byte modulo 256, then $\mathsf{Tuple}\langle\mathsf{Tuple}\langle
\mathsf{u8}^{257}\rangle,\mathsf{u8}\rangle$ and
$\mathsf{Tuple}\langle\mathsf{Tuple}\langle\mathsf{u8}\rangle,\mathsf{u8}^{257}\rangle$ both
encode as $[\mathsf{TUPLE},2][\mathsf{TUPLE},1]$ followed by 258 copies of the encoding of
$\mathsf{u8}$, although they are different types. \Cref{sec:status} reports that the prototype
writes these counts without enforcing the bound, and that the two identity functions on these
types receive the same node hash.
\end{remark}

\subsection{The characterization theorem}

\begin{assumption}[Collision freeness]\label{ass:cr}
For each hash context, distinct preimages that occur have distinct hashes.
\end{assumption}

No hash function is collision free, but a 256-bit function makes collisions improbable, and
derive-key mode keeps contexts apart. The standard idealization quantifies this.

\begin{lemma}[Birthday bound]\label{lem:birthday}
Let $H$ be a uniformly random function into $\{0,1\}^{256}$. For any procedure that makes at
most $q$ distinct queries to $H$, possibly adaptively, the probability that two queries
collide is at most $q(q-1)/2^{257}$.
\end{lemma}
\begin{proof}
The $k$-th distinct query receives a uniformly random value independent of the previous ones,
so it collides with one of the $k-1$ earlier values with probability at most
$(k-1)/2^{256}$. Summing over $k\le q$ gives $q(q-1)/2^{257}$.
\end{proof}

For $q=2^{64}$ preimages the bound is about $2^{-129}$. Modelling each derive-key context of
BLAKE3 as an independent random function is a heuristic~\cite{bellare1993random}; BLAKE3 has a
256-bit output and is designed for a 128-bit security level~\cite{blake3}.
\Cref{ass:cr} is the event that no collision occurs.

\begin{theorem}[Hashes characterize terms up to commutativity]\label{thm:hash-char}
Assume \cref{cond:capacity,ass:cr}. For all operand terms $t,t'$ that arise from well-formed
nodes, $\hash(t)=\hash(t')$ if and only if $t\eqC t'$.
\end{theorem}
\begin{proof}
($\Leftarrow$) Let $t\sim t'$ mean $\hash(t)=\hash(t')$; it is an equivalence relation. It is a
congruence: if $t_k\sim t'_k$ for every $k$, the preimages of $f(t_1,\dots,t_n)$ and
$f(t'_1,\dots,t'_n)$ contain the same operand hashes, sorted the same way, so the op terms have
equal hashes, and so do $\langle f(\dots)\rangle_s$ and $\langle f(\dots)\rangle_s$ for every
$s$. It contains the generators of $\eqC$: for commutative positions, sorting
$(\hash(t_1),\hash(t_2))$ and $(\hash(t_2),\hash(t_1))$ yields the same sequence. Since $\eqC$
is the least congruence containing the generators, ${\eqC}\subseteq{\sim}$.

($\Rightarrow$) By induction on the sum of the sizes of $t$ and $t'$. Suppose
$\hash(t)=\hash(t')$. By \cref{ass:cr}, $\pre(t)=\pre(t')$, and by \cref{lem:prefixfree} $t$ and
$t'$ have the same kind and equal fields. If they are leaves, they are equal. Otherwise
$t=\langle u\rangle_s$ and $t'=\langle u'\rangle_{s}$ with $\Hop(\pre(u))=\Hop(\pre(u'))$, hence
$\pre(u)=\pre(u')$ by \cref{ass:cr}. By \cref{lem:prefixfree}, $u=f(t_1,\dots,t_n)$ and
$u'=f(t'_1,\dots,t'_n)$ have the same op symbol (opcode, mode, attribute and result types)
and the same arity, and their sorted operand-hash sequences are equal. At a non-commutative
position $k$, $\hash(t_k)=\hash(t'_k)$, so $t_k\eqC t'_k$ by induction. At the commutative
positions, the multisets $\{\hash(t_1),\hash(t_2)\}$ and $\{\hash(t'_1),\hash(t'_2)\}$ are
equal; so either $\hash(t_1)=\hash(t'_1)$ and $\hash(t_2)=\hash(t'_2)$, or
$\hash(t_1)=\hash(t'_2)$ and $\hash(t_2)=\hash(t'_1)$. In the first case $t_1\eqC t'_1$ and
$t_2\eqC t'_2$ by induction; in the second, $t_1\eqC t'_2$ and $t_2\eqC t'_1$, and
$f(t_1,t_2,\dots)\eqC f(t'_2,t'_1,\dots)\eqC f(t'_1,t'_2,\dots)$ by the commutativity
generator. Either way $u\eqC u'$, hence $t\eqC t'$.
\end{proof}

The body hash is computed under its own context from the output count, the hashes of the
output operands, and, for each learnable parameter, its type, its bounds as constants and the
hash of its loss node if it has one.

\begin{corollary}[Body hashes]\label{cor:body}
Under the assumptions of \cref{thm:hash-char}, two well-formed nodes have equal body hashes
if and only if they have the same number of outputs, $\unr(a_k)\eqC\unr(a'_k)$ for every output
$k$, and equal learnable boxes.
\end{corollary}

\begin{corollary}[Canonical form]\label{cor:canon}
Under the assumptions of \cref{thm:hash-char}, steps~5 and~6 of canonicalization merge two ops
exactly when their op terms are equal up to commutativity, and the canonical body is
determined by the $\eqC$-classes of the unraveled outputs and of the terms the live invariants
refer to. Consequently canonicalization is idempotent, and each of the following moves of the
specification leaves the canonical body unchanged: consistent renaming of ports, parameters,
bound variables and labels; reordering statements while respecting \ref{wf:order}; duplicating
a pure op; adding a pure op that nothing uses; swapping the operands of a commutative op;
binding a value through an identity op; and changing comments or debug text.
\end{corollary}
\begin{proof}
By \cref{thm:hash-char}, equal op hashes mean $\eqC$-equal op terms, so merging by hash
identifies exactly the $\eqC$-classes; the post-order numbering from the outputs, in port order
and with operands of commutative ops in sorted order, is then a function of the merged graph.
Hence the canonical body depends only on the classes named in the statement, and
canonicalizing it again finds the same classes. Each move preserves those classes: renaming
changes no position, hash or label resolution; reordering does not change $\unr$; a duplicate
of a pure op has the same term; an unused op contributes to no output term; a commutative swap
is a generator of $\eqC$; identity ops are erased before hashing; comments and debug text are
not hashed.
\end{proof}

The specification also states the converse: texts with equal canonical forms are related by
these moves. As listed, the moves do not suffice, because canonicalization makes three further
identifications that the specification states elsewhere: several clauses of one kind join into
one conjunction in source order; a refined parameter or result type is stored as its base type
with the refinement conjoined to the precondition or postcondition; and two literals that
denote the same value, such as \kw{0.99} and its binary32 rounding written out, or the integers
\kw{-0} and \kw{0}, are the same constant. With these three added to the moves, every text can be
moved to the canonical printing of its canonical form (the printer emits callees first and ops in
post-order, which respects \ref{wf:order}), and the moves can be reversed, which gives the
converse. We do not formalize the printer, so we state this direction as a remark rather than a
theorem.

\subsection{Effects are never merged}

Merging ops with equal hashes is sound for pure ops, whose copies and discards are
natural~\cite{fox1976coalgebras}; it would be wrong for effects, since two reads of a clock are
two observations. The specification's answer is the effect token; the next lemma shows that it
suffices.

\begin{lemma}[Effectful ops are never merged]\label{lem:eff-never-merged}
Assume \cref{cond:capacity,ass:cr}, and let $N$ satisfy \ref{wf:eff}. Distinct ops of $N$ that
consume an $\Eff$ value have distinct hashes. Hence step~5 never merges them.
\end{lemma}
\begin{proof}
By \cref{lem:chain}, these ops form the chain $e_1,\dots,e_p$. Let $\epsilon_0$ be input $0$
and $\epsilon_q$ the $\Eff$ result of $e_q$. Suppose $\hash(e_i)=\hash(e_j)$ with $i<j$. By
\cref{ass:cr} and \cref{lem:prefixfree} the two op terms have the same opcode, hence the same
operand layout, and equal operand hashes at every position, in particular at the position of
the $\Eff$ operand: $\hash(\epsilon_{i-1})=\hash(\epsilon_{j-1})$. If $i>1$, both are op results,
so their preimages are equal and $\hash(e_{i-1})=\hash(e_{j-1})$. Repeating $i-1$ times gives
$\hash(\epsilon_0)=\hash(\epsilon_{j-i})$ with $j-i\ge1$. But the preimage of $\epsilon_0$ starts
with byte \texttt{00} (an input) and that of $\epsilon_{j-i}$ with \texttt{01} (an op result),
so the preimages differ and \cref{ass:cr} is contradicted.
\end{proof}

The lemma depends on \ref{wf:eff}. A node in which one token feeds two effect ops cannot be
written in VG-Sym, because the text lowering rejects it, but it can be built through an API or
read from bytes. \Cref{sec:status} reports that the prototype's decoder accepts such a node
and that its merge step then merges two clock reads. The lemma shows that checking
\ref{wf:eff} on every path into the canonicalizer, the decoder included, restores the
guarantee.

\subsection{Equal hashes, equal meaning}

The meaning of an operand term is defined compositionally: $\sem{\mathsf{in}(j,\tau)}=v_j$,
constants denote their values, $\sem{\langle f(t_1,\dots,t_n)\rangle_s}$ applies the meaning of
$f$ (as in \cref{def:eval}) to the meanings of the $t_k$, and the ops on the effect chain read
the trace in the order in which their $\Eff$ operands are nested, which is the order of the
chain.

\begin{lemma}[Unraveling preserves meaning]\label{lem:unr-sem}
For every operand $a$ of a well-formed node, $\val(a)=\sem{\unr(a)}$, both sides being
undefined together.
\end{lemma}
\begin{proof}
By induction on the op index. The two definitions apply the same function to the same operand
values; for the ops on the effect chain, the order of the chain (\cref{lem:chain}) is the order
in which their $\Eff$ operands are nested in the terms, so both definitions hand them the same
trace entries.
\end{proof}

\begin{theorem}[Equal node hashes imply equal meaning]\label{thm:hash-meaning}
Assume \cref{cond:capacity,ass:cr}, and that callees with equal hashes have equal meanings. If
two well-formed canonical nodes $N$ and $N'$ have equal node hashes, then they have the same
format version, effect set, signature and learnable box; for every input, callee semantics and
trace, either both trap or both return the same result; and their contracts are equal up to
the renaming of parameters and bound variables.
\end{theorem}
\begin{proof}
By \cref{ass:cr} the node-hash preimages are equal, so the version, effect set, signature, box,
body hash and contract hash coincide. By \cref{cor:body}, the output terms are $\eqC$-equal. The
meaning of an op symbol in $\mathit{Comm}$ is a commutative function: obviously so for integers
and Booleans, and for floats by \cref{thm:float-comm} below; hence $\eqC$-equal terms have
equal meanings, traps included, by induction on the congruence. By \cref{lem:unr-sem} the
outputs of $N$ and $N'$ have equal values. Ops that only invariants use are covered by the
contract hash, which encodes each port by its operand hash, so their terms, and therefore their
traps, agree as well. Finally, the canonical contract encodes parameters by position and bound
variables by de Bruijn index~\cite{maziarz2021hashing}, so equal contract hashes mean equal
formulas up to renaming.
\end{proof}

Two artifacts with one node hash can still differ outside the node's meaning. The attributes
\kw{export} and \kw{inline} are not hashed, although \kw{export} makes the code generator emit a
checked entry point; the specification leaves that choice to the owner.

\subsection{Why semantic identity needs specifications}

The source document of \VG{} expected two implementations of SHA-256 to converge to one hash.
\Cref{thm:hash-char} shows that a structural hash cannot do this unless the implementations
unravel to the same terms. One might hope for a smarter canonical form. The next theorem shows
that no efficient one exists.

\begin{theorem}[Semantic equivalence is coNP-hard]\label{thm:conp}
Deciding whether two pure, fixed-size, well-formed nodes compute the same function is
coNP-hard, already for nodes with only Boolean inputs and the ops \op{BAND}, \op{BOR} and
\op{BNOT}.
\end{theorem}
\begin{proof}
We reduce the complement of satisfiability, which is coNP-complete~\cite{cook1971complexity},
to equivalence. Given a propositional formula $\psi$ over variables $p_1,\dots,p_n$, build the
node $N_\psi$ with signature $(\mathsf{Bool}^n)\to(\mathsf{Bool})$ that evaluates $\psi$ with
one \op{BAND}, \op{BOR} or \op{BNOT} per connective, and the node $N_\bot$ with the same
signature that returns the constant \kw{false}. Both are well formed and pure, and the
construction takes time linear in $|\psi|$. They compute the same function if and only if
$\psi$ is unsatisfiable.
\end{proof}

\begin{corollary}\label{cor:no-canon}
If some polynomial-time computable function $c$ on nodes satisfied $c(N)=c(N')$ exactly when
$N$ and $N'$ compute the same function, then $\mathrm{P}=\mathrm{NP}$.
\end{corollary}
\begin{proof}
Equivalence would be decidable in polynomial time by comparing $c(N_\psi)$ with $c(N_\bot)$,
so unsatisfiability would be in P, hence coNP $=$ P and therefore NP $=$ P.
\end{proof}

Any efficiently computable identity must therefore either merge some inequivalent programs or
separate some equivalent ones. \VG{} chooses the second, which is the safe failure, and handles
semantic identity by specification: a SPEC node is an executable reference, and an
implementation joins its spec class when its postcondition \kw{result == spec:s(\dots)} is
proven. The registry records how membership was shown: \emph{proved} for every input,
\emph{bounded}$(n)$ for inputs of size at most $n$, or \emph{tested}$(m)$ on $m$ cases, which a
caller's proof may never rely on.

\begin{proposition}[Spec classes are extensional]\label{prop:specclass}
Let $I_1$ and $I_2$ satisfy their contracts (\cref{def:satisfies}) with postconditions
$\kw{result == spec:}s(\bar x)$, for a result type on which \kw{==} is identity (integers,
Booleans and collections of them), and let $s$ satisfy its own contract. Then $I_1(\bar v)=
I_2(\bar v)$ for every $\bar v$ that satisfies the preconditions of $I_1$, $I_2$ and $s$.
\end{proposition}
\begin{proof}
For such $\bar v$, $\bar v$ lies in the domain of $\kw{spec:}s$, where the term denotes
$s(\bar v)$ in every interpretation. The postconditions of $I_1$ and $I_2$ hold at $\bar v$, so
$I_1(\bar v)=s(\bar v)=I_2(\bar v)$.
\end{proof}

For float results \kw{==} is IEEE equality, under which $-0=+0$ and $\NaN\neq\NaN$; a spec class
over floats must use \kw{same} to obtain the same conclusion.

\subsection{Short hashes}

Text and diagnostics show a hash as \texttt{\#} followed by 12 hexadecimal digits. If $n$
uniformly random 48-bit prefixes are drawn, a collision becomes more likely than not near
$n\approx\sqrt{2\ln 2\cdot 2^{48}}\approx 1.98\times10^{7}$; for 8 digits (32 bits) the same
estimate is about $7.7\times10^{4}$. The toolchain lengthens a displayed prefix on any clash
within a module, and binaries always carry all 32 bytes, so short hashes are a convenience,
never an identity.
